\chapter{2024年上海卷（秋）}

\section{填空题}

\begin{problem}
设全集 \(U=\{1,2,3,4,5\}\)，集合 \(A=\{2,4\}\)，则 \(\bar{A}=\)\fillinblank{}.
\end{problem}

\begin{problem}
已知 \(f(x)=\begin{cases}
\sqrt{x}, & x>0,\\
1, & x\le 0,
\end{cases}\) 则 \(f(3)=\)\fillinblank{}.
\end{problem}

\begin{problem}
已知 \(x\in \R\)，则不等式 \(x^2-2x-3<0\) 的解集为\fillinblank{}.
\end{problem}

\begin{problem}
已知 \(f(x)=x^3+a\)，\(x\in \R\)，且 \(f(x)\) 是奇函数，则 \(a=\)\fillinblank{}.
\end{problem}

\begin{problem}
已知 \(k\in \R\)，\(\bs{a}=(2,5)\)，\(\bs{b}=(6,k)\)，且 \(\bs{a}\parallel\bs{b}\)，则 \(k=\)\fillinblank{}.
\end{problem}

\begin{problem}
在 \((x+1)^n\) 的二项展开式中，若各项系数和为 32，则 \(x^2\) 项的系数为\fillinblank{}.
\end{problem}

\begin{problem}
已知抛物线 \(y^2=4x\) 上有一点 \(P\) 到准线的距离为 9，那么点 \(P\) 到 \(x\) 轴的距离为\fillinblank{}.
\end{problem}

\begin{problem}
小王参加知识竞赛，题库中 A 组题有 \(5000\) 道，B 组题有 \(4000\) 道，C 组题有 \(3000\) 道，若小王做对这三组题的概率依次为 \(0.92\)，\(0.86\)，\(0.72\)，则随机从题库中抽取一道题，小王做对的概率是\fillinblank{}.
\end{problem}

\begin{problem}
设 \(m\in \R\)，若虚数 \(z\) 的实部为 1，且 \(z+\frac{2}{z}=m\)，则 \(m=\)\fillinblank{}.
\end{problem}

\begin{problem}
已知某集合中的元素是不重复的数字组成的三位正整数，若该集合中任意两个数的积均为偶数，则该集合的元素个数的最大值为\fillinblank{}.
\end{problem}

\begin{problem}
海上有灯塔 \(O\)，\(A\)，\(B\) 和船只 \(T\)，\(A\) 在 \(O\) 的正东方向，\(B\) 在 \(O\) 的正北方向，\(OA=OB\)，\(O\)，\(A\)，\(T\)，\(B\) 按逆时针排列. 若 \(\angle ATO=37^\circ\)，\(\angle BTO=16.5^\circ\)，则 \(\angle BOT=\)\fillinblank{}（精确到 \(0.1^\circ\)）.

\bitmapfigure[max width=0.42\linewidth]{img/2024/shanghai/q11_fig1.png}
\end{problem}

\begin{problem}
等比数列 \(\{a_n\}\) 满足首项 \(a_1>0\)，公比 \(q>1\)，记 \(I_n=\{x-y\mid x,y\in [a_1,a_2]\cup [a_n,a_{n+1}]\}\)，若对任意正整数 \(n\)，集合 \(I_n\) 都是闭区间，则 \(q\) 的取值范围是\fillinblank{}.
\end{problem}

\section{单选题}

\begin{problem}
已知气温（单位 \(^{\circ}\mathrm C\)）和海水表层温度（单位：摄氏度）相关，且相关系数为正数，对此描述正确的是（\quad）

\choices
    {随着气温由低变高，海水表层温度由低变高}
    {随着气温由低变高，海水表层温度由高变低}
    {随着气温由低变高，海水表层温度有由低变高的趋势}
    {随着气温由低变高，海水表层温度有由高变低的趋势}
\end{problem}

\begin{problem}
下列函数 \(f(x)\) 的最小正周期是 \(2\pi\) 的是（\quad）

\choices
    {\(y=\sin x+\cos x\)}
    {\(y=\sin x\cos x\)}
    {\(y=\sin^2 x+\cos^2 x\)}
    {\(y=\sin^2 x-\cos^2 x\)}
\end{problem}

\begin{problem}
已知空间直角坐标系 \(Oxyz\) 中的点集 \(\Omega\)，对任意的 \(P_1\)，\(P_2\)，\(P_3\in \Omega\)，均存在不全为 \(0\) 的实数 \(\lambda_1\)，\(\lambda_2\)，\(\lambda_3\)，使得 \(\lambda_1\overrightarrow{OP_1}+\lambda_2\overrightarrow{OP_2}+\lambda_3\overrightarrow{OP_3}=\bs{0}\). 若 \((1,0,0)\in \Omega\)，则 \((0,0,1)\notin \Omega\) 的一个充分条件是（\quad）

\choices
    {\((0,0,0)\in \Omega\)}
    {\((-1,0,0)\in \Omega\)}
    {\((0,-1,0)\in \Omega\)}
    {\((0,0,-1)\in \Omega\)}
\end{problem}

\begin{problem}
已知函数 \(f(x)\) 的定义域为 \(\R\)，定义集合 \(M=\{x_0\mid x_0\in \R,\ \forall x\in(-\infty,x_0),\ f(x)<f(x_0)\}\). 在使得 \(M=[-1,1]\) 的所有 \(f(x)\) 中，下列成立的是（\quad）

\choices
    {存在 \(y=f(x)\) 是偶函数}
    {存在 \(y=f(x)\) 在 \(x=2\) 处取最大值}
    {存在 \(y=f(x)\) 是严格增函数}
    {存在 \(y=f(x)\) 在 \(x=-1\) 处取到极小值}
\end{problem}

\section{解答题}

\begin{problem}
如图，在正四棱锥 \(P-ABCD\) 中，\(O\) 为底面 \(ABCD\) 的中心.

\begin{enumerate}
\item 若 \(AP=5\)，\(AD=3\sqrt{2}\)，将 \(\triangle POA\) 绕直线 \(PO\) 旋转一周，求所得旋转体的体积；
\item 若 \(AP=AD\)，\(E\) 为 \(PB\) 的中点，求直线 \(BD\) 与平面 \(AEC\) 所成角的大小.
\end{enumerate}

\bitmapfigure[max width=0.55\linewidth]{img/2024/shanghai/q17_fig1.png}
\end{problem}

\begin{problem}
若 \(f(x)=\log_a x\)（\(a>0\)，\(a\ne 1\)）.

\begin{enumerate}
\item \(y=f(x)\) 过 \((4,2)\)，求 \(f(2x-2)<f(x)\) 的解集；
\item 存在 \(x\) 使得 \(f(x+1)\)，\(f(ax)\)，\(f(x+2)\) 成等差数列，求 \(a\) 的取值范围.
\end{enumerate}
\end{problem}

\begin{problem}
某地区为调查初中学生体育锻炼时长与学业成绩的关系，从该地区 29000 名初中生中抽取 580 人，得到日均体育锻炼时长（表中简称时长）与学业成绩的数据如下表所示：

\begin{center}
\begin{tabular}{|c|c|c|c|c|c|c|}
\hline
时长 & \([0,0.5)\) & \([0.5,1)\) & \([1,1.5)\) & \([1.5,2)\) & \([2,2.5)\) & 合计 \\
\hline
优秀 & 5 & 44 & 42 & 3 & 1 & 95 \\
\hline
不优秀 & 134 & 147 & 137 & 40 & 27 & 485 \\
\hline
合计 & 139 & 191 & 179 & 43 & 28 & 580 \\
\hline
\end{tabular}
\end{center}

\begin{enumerate}
\item 估计该地区 29000 名初中生中体育锻炼时长大于 1 小时人数；
\item 估计该地区初中生平均日均体育锻炼时长（结果精确到 0.1 小时）；
\item 判断是否有 95\% 的把握认为该地区初中生学业成绩优秀与日均体育锻炼时长不小于 1 小时且小于 2 小时有关？
（附：\(\chi^2=\frac{n(ad-bc)^2}{(a+b)(c+d)(a+c)(b+d)}\)，其中 \(n=a+b+c+d\)，\(P(\chi^2\ge 3.841)\approx 0.05\).）
\end{enumerate}
\end{problem}

\begin{problem}
已知双曲线 \(\Gamma:x^2-\frac{y^2}{b^2}=1\)（\(b>0\)），左右顶点分别为 \(A_1\)，\(A_2\)，过点 \(M(-2,0)\) 的直线 \(l\) 交双曲线 \(\Gamma\) 于 \(P\)，\(Q\) 两点.

\begin{enumerate}
\item 若离心率 \(e=2\)，求 \(b\) 的值；
\item 若 \(b=\frac{2\sqrt{6}}{3}\)，点 \(P\) 在第一象限，\(\triangle MA_2P\) 为等腰三角形，求点 \(P\) 的坐标；
\item 连接 \(QO\)（\(O\) 为坐标原点）并延长交双曲线 \(\Gamma\) 于点 \(R\)，若 \(\overrightarrow{A_1R}\cdot\overrightarrow{A_2P}=1\)，求 \(b\) 的取值范围.
\end{enumerate}
\end{problem}

\begin{problem}
设 \(D\) 是 \(\R\) 的一个非空子集，\(y=f(x)\) 是定义在 \(D\) 上的函数. 对于点 \(M(a,b)\)，记 \(s(x)=(x-a)^2+\left( f(x)-b \right)^2\). 若对于点 \(P(x_0,f(x_0))\)，函数 \(y=s(x)\) 在 \(x=x_0\) 处取得最小值，则称 \(P\) 是 \(M\) 的 \(f\) 最近点.

\begin{enumerate}
\item 对于 \(f(x)=\frac{1}{x}\)，\(M(0,0)\)，\(D=(0,+\infty)\)，求证：存在 \(M\) 的 \(f\) 最近点；
\item 对于 \(f(x)=\e^x\)，\(M(1,0)\)，\(D=\R\)，若曲线 \(y=f(x)\) 上一点 \(P\) 满足 \(MP\) 垂直于 \(y=f(x)\) 在点 \(P(x_0,f(x_0))\) 处的切线，则 \(P\) 是否为 \(M\) 的 \(f\) 最近点？
\item 已知 \(D=\R\)，\(y=f(x)\) 的导函数为 \(y=f'(x)\)，且 \(y=g(x)\) 在 \(\R\) 上的函数值恒正. 对任意 \(t\in \R\)，点 \(M_1\) 的坐标为 \((t-1,f(t)-g(t))\)，点 \(M_2\) 的坐标为 \((t+1,f(t)+g(t))\). 若对任意的 \(t\in \R\)，总存在 \(y=f(x)\) 上一点 \(P\)，使得 \(P\) 既是 \(M_1\) 的 \(f\) 最近点，又是 \(M_2\) 的 \(f\) 最近点，试求 \(y=f(x)\) 的单调性.
\end{enumerate}
\end{problem}

